\textbf{Selección vigente CEN-41.} Se sustituyen los dos LDO negativos TPS7A3301 por dos positivos adicionales, total cuatro TPS7A4700 en cuatro celdas flotantes. Se seleccionan trim OEM1,72 kohm, cuatro LTC4365 y ocho MOS de corte. Las conexiones y pasivos negativos que siguen pertenecen al antecedente sustituido, sin otra compra. Se conservan cuatro interfaces GCI-38 y cuatro TEN/PTCB; secuencia bipolar, ventana válida, histéresis y respuesta dinámica siguen pendientes. El presupuesto DC global vigente con ENV-40 es356,7552 W; las pérdidas posteriores a TEN no se agregan otra vez.

\textbf{Alimentación diagnóstica vigente CEN/ECI-38.} Se conserva la selección LEM, burdens y AI de este desarrollo, exclusivamente operacional. La salida simple nominal de los TEN es 15 V; el diseño vigente selecciona objetivo de ajuste de prerregulación a 16,5 V y dos LDO por banco para salida cercana a ±14,9 V. No se conecta TEN directamente a LEM. Trim OEM, protección independiente, secuencia/arranque y PCB siguen pendientes. Las cuatro interfaces de estado elegidas en GCI-38 sustituyen la reserva genérica de interfaz, sin habilitar recarga ni cerrar fallos. El subtotal global control/energía pasa de la base 348,7552 W de esta memoria a 352,7552 W DC, con 4 W propios adicionales excluyendo DI ya contadas; no se duplican pérdidas LDO dentro de TEN.

\section{Medida operacional por cadena de batería (CEN-35)}
\label{sec:sd4-medida-cen35}
Se seleccionan cuatro transductores LEM IT 200-S ULTRASTAB, referencia 88.36.44.000.0: dos por banco, uno en el conductor de cada cadena. Su función es registrar corriente bipolar, reparto y suma positiva de recarga; no sustituyen la protección independiente de gas/caudal. La ficha versión 4 publica rango DC continuo y de medida $\pm200$ A, salida $\pm200$ mA y relación 1:1000. Se orienta la flecha de cada primario hacia la batería: carga positiva, descarga negativa. La suma es $I_{\rm carga}=\sum_{i=1}^{2}\max(0,I_i)$ por banco; una cadena descargando no compensa la emisión de otra cargando. En un cribado propio, 60 kW a 400 V DC con eficiencia 0,94 requieren 159,5745 A si una sola cadena toma toda la descarga; quedan dentro del rango de medida, pero ese cribado no acredita reparto, protección del conductor ni potencia completa de la instalación \parencite[pp.~1, 5 y 10]{lote35-lem200}.

La ficha excluye usos de soporte vital o críticos para la seguridad y aplicaciones donde un fallo pueda producir lesiones o daños mayores. Por ello la selección se limita a diagnóstico operacional: no se le atribuye SIL/PL ni habilitación de recarga por este sensor. Gas, caudal y retirada independiente conservan sus funciones y pendientes de INT-31. Se conserva el límite requerido de 20 A reales agregados por banco de CEN-32; aún no hay rango, paso, tolerancia ni interfaz OEM del cargador 93T que permitan habilitar ese perfil. Registrar 20 A no constituye limitar el cargador a 20 A \parencite[p.~11]{lote35-lem200}.

\textbf{Instalación y aislamiento.} Se fija montaje recto con dos tornillos M5 por sensor, par 3,7 Nm; ocho tornillos en total. El paso primario admite diámetro máximo 26 mm. Se instala dentro de gabinete cerrado que impida acceso a partes peligrosas, con separación mínima 30 mm respecto a terminales primarios/componentes vecinos; ambiente de operación 10--50 $^\circ$C, humedad 20--80\,\% sin condensación y conductor primario a máximo 50 $^\circ$C. Se proyectan cables secundarios aptos para 100 $^\circ$C, sujetos y separados del primario. La coordinación publicada es reforzada 1.000 V CAT II/PD2 según IEC 61010-1 y 600 V CAT III/PD2 según EN 50178. El máximo proyectado 607 V DC excede los 600 V del segundo ejemplo; el ensayo 5,4 kV AC no equivale a tensión de trabajo. Antes de liberar se debe acreditar categoría de sobretensión aplicable al gabinete, cable primario real y aislamiento complementario cuando proceda. La tabla de cables HAR07 sólo publica 2.350 V CAT III para aislamiento simple; no se la transforma en reforzado. La selección operacional no cierra esta coordinación \parencite[pp.~2--4 y 10]{lote35-lem200}.

\textbf{Conversión y canales seleccionados.} Se seleccionan cuatro resistencias de cuatro terminales Kelvin VPG VPR221Z, 25 $\Omega$, tolerancia nominal $\pm0,01\,\%$, referencia Y169025R0000T9L. Cada resistencia se coloca localmente entre OUTPUT pin 6 y RETURN pin 1 del D-SUB-9 del sensor; pares Kelvin de medida hacia el módulo, sin usar el cable como parte del burden. A 200 mA entrega 5 V y disipa 1 W; a 20 A primarios entrega 0,5 V. Los 25 $\Omega$ se encuentran dentro del límite de 30 $\Omega$ de la curva OEM. La familia publica 1,5 W al aire a 25 $^\circ$C y 8 W sobre disipador; no se atribuyen 8 W al aire. Se fija montaje térmico sobre placa disipadora y comprobación de temperatura/resistencia instalada. El valor máximo TCR de esta familia/rango es $\pm1,8$ ppm/K; los datos típicos no se usan como máximos. La copia disponible de la ficha de familia es revisión 25-03-2010: se cotejará la revisión vigente del pedido, sin presentar ese cotejo como realizado \parencites[pp.~8 y 10]{lote35-lem200}[pp.~1--3]{lote35-vprfamily}{lote35-vprsku}.

Se seleccionan dos WAGO 750-479, uno por banco: dos entradas diferenciales $\pm10$ V por módulo, cuatro canales efectivos para cuatro cadenas. Son adicionales al módulo 4--20 mA de ENV-34. La escala es $V=0,025I$ y $I=40V$. El módulo publica impedancia de entrada 1 M$\Omega$, LSB 1,2 mV, desviación máxima 0,05\,\% del límite superior de rango a 25 $^\circ$C y 0,4\,\% en todo el intervalo térmico. La separación 500 V del módulo no resuelve aislamiento primario de 607 V: corresponde al lado secundario de medida. Cada módulo consume 100 mA a 5 V del bus; ENV-34 incorpora 432 mA por nodo frente a capacidad 1.700 mA. Las cuatro señales de estado LEM ocupan dos DI por banco; no se declara supervisión de rotura de cable analógico a cero sólo por tener un canal \parencite[pp.~1--3]{lote35-wago479}.

\textbf{Estado y conexión.} En la figura OEM del D-SUB-9, pin 9 es $+U_C$, pin 5 es $-U_C$, pin 6 OUTPUT, pin 1 RETURN y pin 4 Ground/pantalla; pines 2 y 7 NC. El estado usa colector pin 8 y emisor pin 3: transistor ON indica alimentaciones presentes, detector de cero operativo y compensación no superior al 110\,\% nominal. OFF invalida la lectura. El estado admite alimentación 4--45 V y corriente 2--30 mA. La interfaz hacia DI 24 V deberá acreditar polaridad, corriente, umbral y detección de cable cortado; no se fija una resistencia genérica como interfaz ya resuelta ni se conecta directamente a A1 de XPS, cuya corriente nominal supera 30 mA. Un corto que mantenga ON o una salida de medida congelada requiere diagnóstico adicional. En sobrecarga la salida barre entre $\pm200$ mA y el estado abre; recuperación automática puede tardar segundos. Ninguna lectura próxima a cero durante ese barrido cuenta como carga nula. El retardo de sensor 1 $\mu$s corresponde al ensayo de pendiente 100 A/$\mu$s; no es tiempo completo sensor--PLC--actuador. El módulo publica demora/repetición 1 ms, pero faltan tiempo de tarea, comunicaciones, lógica y actuación \parencites[pp.~5, 7--8 y 10]{lote35-lem200}[pp.~2--3]{lote35-wago479}.

\textbf{Presupuesto parcial de error.} A rango térmico completo, por canal WAGO se usan $0,004(10)/0,025=1,6$ A y un LSB completo $0,0012/0,025=0,048$ A. Para LEM se suman offset $80$ ppm de 200 A, deriva $2$ ppm/K durante 25 K, estabilidad $1$ ppm/mes durante doce meses y linealidad $3$ ppm: 0,029 A por canal. Para dos cadenas el término aditivo es $2(1,6+0,048+0,029)=3,354$ A. Se fija reserva propia de ganancia 0,1\,\% para burden instalado y 25 ppm adicionales por carga de entrada. Bajo esas hipótesis, $I_{\rm real}^{+}\le(I_{\rm medido}^{+}+3,354)/(1-0,001025)$: una suma medida de 20 A puede representar 23,3780 A; una de 17 A, 20,3749 A. Este segundo valor deja sólo 0,0942 A respecto al máximo real de cribado 20,4691 A de INT-31. Se fija alarma operacional a 17 A, sin autorizar recarga. El presupuesto no incluye un máximo completo verificado de error de relación, interferencia/rizado, montaje ni todos los fallos; además, LEM define algunos máximos como intervalo estadístico 99,73\,\%, no garantía universal. No se declara cerrado el límite de gas con esta suma parcial ni se extrapola ruido RMS a cota máxima \parencite[pp.~5--6]{lote35-lem200}.

\section{Alimentación de medida y conciliación de consumos (ECI/GCI-35)}
\label{sec:sd4-alimentacion-eci35}
Se seleccionan cuatro TRACO TEN 30-2413, entrada 18--36 V DC, salida simple aislada 15 V/2 A: dos por banco, alimentados desde su QUINT 24 V. En cada banco se unen el negativo de la salida positiva y el positivo de la salida negativa para formar referencia 0; los extremos alimentan $+15$ y $-15$ V de sus dos LEM. No se unen referencias de bancos A/B ni se utiliza salida dual con regulación cruzada del 5\,\%. Cada sensor exige $\pm14,25$--$\pm15,75$ V. La suma estática de exactitud 1\,\%, línea 0,2\,\%, carga 0,5\,\% y temperatura 0,02\,\%/K a 25 K da 2,2\,\%: banda calculada 14,67--15,33 V. No cubre rizado/transitorios ni fallo: rizado 150 mVpp es típico y protección de sobretensión 125\,\% típica puede superar el máximo absoluto LEM 15,75 V. Por ello monitor/limitación de ambas alimentaciones, fallo de referencia y respuesta de arranque siguen pendientes y conservan la medida sin habilitación de seguridad \parencites[pp.~1--4]{lote35-ten30}[pp.~3 y 5]{lote35-lem200}.

Se seleccionan cuatro PTCB E1 24DC/2A NO, referencia 2909903, adicionales a los dos del SEL; cada convertidor tiene rama propia. Se reserva por diseño 2 A a 24 V por convertidor: 48 W cada uno, 192 W globales. Este valor es reserva de proyecto, no consumo OEM máximo ni eficiencia mínima: la ficha publica 1,5 A a plena carga y eficiencia 91\,\%, ambos típicos para la variante. Con 0,8 W por PTCB, el incremento reservado es 195,2 W, 97,6 W por QUINT. Los dos LEM de cada banco, calculados conservadoramente como $2(15)(0,080+0,200)$ por unidad, quedan en 16,8 W a rango completo; ese consumo y los burdens ya están aguas abajo de los convertidores y no se suman de nuevo. Tampoco se suman 0,5 W de cada AI como carga 24 V independiente del bus CPU de ENV-34 \parencites[pp.~1--3]{lote35-ten30}[p.~5]{lote35-lem200}{lote29-ptcb2}.

La reserva anterior exclusivamente energética de 72,8 W pasa a subtotal 268 W, 134 W por QUINT, antes de control ambiental. Partiendo del subtotal integrado INT-31 de 101,6 W, el resultado es 296,8 W globales, 148,4 W por QUINT, antes de las adiciones de ENV-34. Estas bases son diferentes y no se suman entre sí. La reserva AC ya establecida para las fuentes QUINT no vuelve a recibir esos 195,2 W DC como una segunda carga AC; sí se debe verificar capacidad, pérdidas y arranque de las fuentes reales. TRACO recomienda fusible de entrada 3,15 A lento para variantes 24 V: el PTCB 2 A seleccionado no acredita esa curva ni paso de arranque y no se declara equivalente. Protecciones/cableado definitivos y montaje de convertidores sobre placa quedan sujetos a coordinación y ensayo de carga/arranque. No se aumenta autonomía ni potencia PRP del grupo por estos subtotales. GCI-32 conserva medida AC, permiso y pendientes de mapa de pedido, tiempos, fallas/coord y transitorios; ECI-32 conserva balance completo por fase/estado y perfil anual pendientes.

La conciliación vigente con ENV-34 alcanza 348,7552 W DC globales y 174,3776 W por QUINT, antes de receptores pendientes; los subtotales anteriores no se suman de nuevo. El margen nominal 65,6224 W por fuente no acredita arranque ni autonomía.
